\documentclass[aspectratio=169,11pt]{ctexbeamer}

\usetheme{Madrid}
\usecolortheme{default}
\setbeamertemplate{navigation symbols}{}
\setbeamertemplate{footline}[frame number]

\usepackage{amsmath,amssymb,bm}
\usepackage{booktabs}
\usepackage{tikz}
\usetikzlibrary{arrows.meta,positioning}

\title[rydcalc 与寿命计算]{rydcalc 功能简介与 Rydberg 态寿命计算}
\subtitle{重点：\texttt{total\_decay} 如何从态、环境和偶极矩阵元得到寿命}
\author{Noise-Tolerant Quantum Control Optimization}
\date{\today}

\newcommand{\ket}[1]{\left|#1\right\rangle}
\newcommand{\bra}[1]{\left\langle#1\right|}
\newcommand{\Hz}{\mathrm{Hz}}
\newcommand{\GHz}{\mathrm{GHz}}
\newcommand{\au}{\mathrm{a.u.}}

\begin{document}

\begin{frame}
  \titlepage
\end{frame}

\begin{frame}{这份 slides 讲什么}
  \begin{itemize}
    \item \texttt{rydcalc} 是项目中用于 Rydberg 原子结构、矩阵元、Stark 图、pair interaction 和寿命估算的外部子模块。
    \item 本项目主要通过 \texttt{neutral\_yb.external.rydcalc\_adapter} 调用它，避免直接改子模块。
    \item 本报告重点解释：
      \begin{itemize}
        \item \texttt{rydcalc} 能提供哪些原子物理量；
        \item \texttt{total\_decay(st, env)} 如何自动构造终态并求和；
        \item \texttt{partial\_decay} 中黑体辐射和自发辐射因子如何进入；
        \item 对 \({}^{171}\mathrm{Yb}\) MQDT 态使用时需要注意哪些数值细节。
      \end{itemize}
  \end{itemize}
\end{frame}

\begin{frame}{rydcalc 在项目中的位置}
  \begin{columns}[T]
    \begin{column}{0.52\textwidth}
      \begin{itemize}
        \item 子模块位置：\texttt{rydcalc/}
        \item 项目适配层：\texttt{rydcalc\_adapter.py}
        \item 当前 notebook 示例：
          \begin{itemize}
            \item UV Rabi / power estimate notebook；
            \item \(C_3/C_6\) estimate notebook；
            \item 65S lifetime diagnostic notebook；
            \item windowed \texttt{get\_nearby} repair notebook；
            \item hybrid + residual recalculation notebook。
          \end{itemize}
      \end{itemize}
    \end{column}
    \begin{column}{0.43\textwidth}
      \begin{block}{适配层的主要职责}
        \begin{itemize}
          \item 检查 \texttt{rydcalc} 子模块存在；
          \item 检查 ARC C extension 是否为当前 Python 构建；
          \item 为 NumPy 2 / SciPy API 差异提供兼容 shim；
          \item 提供 \texttt{build\_yb171\_atom()} 入口。
        \end{itemize}
      \end{block}
    \end{column}
  \end{columns}
\end{frame}

\begin{frame}{rydcalc 能做什么}
  \begin{table}
    \centering
    \small
    \begin{tabular}{lll}
      \toprule
      功能 & 典型 API & 本项目用途 \\
      \midrule
      原子对象 & \texttt{Ytterbium171}, \texttt{Rubidium} & 构造物种和能级模型 \\
      态构造 & \texttt{atom.get\_state(...)} & 生成 Rydberg / NIST 态对象 \\
      矩阵元 & \texttt{get\_multipole\_me} & UV Rabi、跃迁强度、衰变率 \\
      Stark 图 & \texttt{analysis\_stark} & 电场混合、能级位移 \\
      Pair interaction & \texttt{analysis\_pair\_interaction} & \(C_3,C_6\)、blockade 估计 \\
      寿命 & \texttt{total\_decay}, \texttt{partial\_decay} & Rydberg decay / blackbody rate \\
      \bottomrule
    \end{tabular}
  \end{table}

  \vspace{0.5em}
  对 \({}^{171}\mathrm{Yb}\) 而言，核心价值是 MQDT 模型：Rydberg 态不是单一氢样态，而是多个 channel 的线性组合。
\end{frame}

\begin{frame}{基本调用流程}
  \begin{block}{典型代码}
    \small
    \texttt{rydcalc = load\_rydcalc(require\_c\_extension=True)}\\
    \texttt{yb = build\_yb171\_atom(cpp\_numerov=True, use\_db=False)}\\
    \texttt{st = yb.get\_state((65, 0, 1.5, -1.5), tt="vlfm")}\\
    \texttt{env = rydcalc.environment(T\_K=300)}\\
    \texttt{tau = 1 / yb.total\_decay(st, env)}
  \end{block}

  \begin{itemize}
    \item \texttt{tt="vlfm"} 对 Yb171 Rydberg MQDT 态表示 \((\nu/L/F/m_F)\) 输入格式。
    \item \texttt{environment} 包含温度、电磁场、LDOS 等外部环境参数。
    \item \texttt{total\_decay} 返回总衰变率 \(\Gamma\)，寿命是 \(\tau = 1/\Gamma\)。
  \end{itemize}
\end{frame}

\begin{frame}{态对象里保存了什么}
  \begin{itemize}
    \item 能量：\texttt{state.get\_energy\_Hz()}。
    \item 角动量量子数：例如 \(L,F,m_F\) 或 \(n,l,j,m\)。
    \item MQDT channel 列表：\texttt{state.channels}。
    \item channel 权重：\texttt{state.Ai} 和相关 MQDT 系数。
    \item 径向波函数信息：需要时由 Numerov 或 Whittaker 方法生成。
  \end{itemize}

  \begin{block}{对寿命计算最关键的接口}
    \[
      \texttt{state.get\_multipole\_me(other, k=1)}
      \quad\Rightarrow\quad
      \bra{\mathrm{other}} er_q \ket{\mathrm{state}}
    \]
    这里返回的是原子单位下的电偶极矩阵元；\texttt{partial\_decay} 会再转成 SI 单位进入速率公式。
  \end{block}
\end{frame}

\begin{frame}{environment 对寿命的作用}
  \texttt{rydcalc.environment} 当前保存：
  \begin{itemize}
    \item \texttt{T\_K}：温度，决定黑体光子占据数；
    \item \texttt{Bz\_Gauss}, \texttt{Ez\_Vcm}：场参数，主要给 Stark / Zeeman 分析使用；
    \item \texttt{ldos(f, q)}：局域光学态密度修正，默认恒等于 1；
    \item \texttt{diamagnetism}：是否包含抗磁项。
  \end{itemize}

  \begin{block}{寿命计算中真正用到的两项}
    \[
      T \Rightarrow n_{\mathrm{th}}(\omega,T),
      \qquad
      \mathrm{LDOS}(f,q) \Rightarrow \mathrm{Purcell/geometry\ correction}.
    \]
  \end{block}
\end{frame}

\begin{frame}{total\_decay 的总体结构}
  \begin{columns}[T]
    \begin{column}{0.49\textwidth}
      \begin{enumerate}
        \item 初始化总速率：
          \[
            \Gamma_{\mathrm{tot}}=0.
          \]
        \item 构造一个 \texttt{single\_basis()}。
        \item 自动填充可能衰变到的终态。
        \item 对每个终态 \(s_f\) 调用：
          \[
            \Gamma_{i\rightarrow f}
            =
            \texttt{partial\_decay}(s_i,s_f,\mathrm{env}).
          \]
        \item 累加：
          \[
            \Gamma_{\mathrm{tot}}=\sum_f\Gamma_{i\rightarrow f}.
          \]
      \end{enumerate}
    \end{column}
    \begin{column}{0.46\textwidth}
      \begin{block}{返回值}
        \begin{itemize}
          \item \texttt{printstates=False}：返回 \(\Gamma_{\mathrm{tot}}\)。
          \item \texttt{printstates=True}：返回非零 partial decay 列表。
          \item 寿命不直接返回，要手动取倒数：
            \[
              \tau = 1/\Gamma_{\mathrm{tot}}.
            \]
        \end{itemize}
      \end{block}
    \end{column}
  \end{columns}
\end{frame}

\begin{frame}[fragile]{total\_decay 的关键代码逻辑}
  \scriptsize
  \begin{verbatim}
gamma_tot = 0
sb = single_basis()

nmax = st.n + 10
sb.fill(st, {
    'dn': st.n,
    'dl': 1,
    'dm': 1,
    'include_fn': lambda s, s0: True if s.n < nmax else False,
})

for s in sb.states:
    pd = self.partial_decay(st, s, env)
    gamma_tot += pd

return gamma_tot
  \end{verbatim}

  \vspace{-0.5em}
  这段逻辑说明：\texttt{total\_decay} 的物理核心不在传播动力学，而在“列出终态 + 对所有 E1 partial rate 求和”。
\end{frame}

\begin{frame}{自动终态生成的调用链}
  \begin{itemize}
    \item \texttt{total\_decay} 本身不直接知道有哪些终态。
    \item 它先构造 \texttt{single\_basis()}，再调用：
      \[
        \texttt{single\_basis.fill(st, include\_opts)}.
      \]
    \item \texttt{fill} 再调用初态所属 atom 的：
      \[
        \texttt{st.atom.get\_nearby(st, include\_opts)}.
      \]
    \item 对氢样 / 普通碱金属态，\texttt{get\_nearby} 的直觉是按
      \(\Delta n,\Delta l,\Delta m\) 枚举附近主量子数。
    \item 对 \({}^{171}\mathrm{Yb}\) MQDT 态，实际逻辑变成按有效量子数 \(\nu\) 做 MQDT 求根。
  \end{itemize}
\end{frame}

\begin{frame}[fragile]{Yb171 \texttt{get\_nearby} 的关键代码}
  \tiny
  \begin{verbatim}
o = {'dn': 2, 'dl': 2, 'dm': 1, 'ds': 0}
for k, v in include_opts.items():
    o[k] = v
if 'df' not in o.keys():
    o['df'] = o['dl']

nu0 = st.nub

for l in np.arange(st.channels[0].l - o['dl'],
                   st.channels[0].l + o['dl'] + 1):
    if l < 0:
        continue
    for f in np.arange(st.f - o['df'], st.f + o['df'] + 1):
        if f < 0:
            continue
        solver = [d for d in self.mqdt_models
                  if d['L'] == l and d['F'] == f][0]['model']

        boundstatesinrange = solver.boundstatesinrange([
            nu0 - o['dn'],
            nu0 + o['dn'],
        ])
  \end{verbatim}

  \vspace{-0.7em}
  这段代码说明：\texttt{dn} 在 Yb171 路径中直接成为 \(\nu\) 的搜索半宽；
  \texttt{include\_fn} 不参与这里的 \texttt{boundstatesinrange} 区间构造。
\end{frame}

\begin{frame}{Yb171 MQDT 终态如何生成}
  对 \({}^{171}\mathrm{Yb}\) MQDT 态，\texttt{get\_nearby} 会：
  \begin{itemize}
    \item 从初态 channel 读出 \(L_i\)，从初态对象读出 \(F_i,m_F\)；
    \item 按输入的 \texttt{dl}, \texttt{df}, \texttt{dm} 扫描：
      \[
        L=L_i-\Delta L,\ldots,L_i+\Delta L,
        \qquad
        F=F_i-\Delta F,\ldots,F_i+\Delta F.
      \]
    \item 对每个 \((L,F)\) 选择对应 MQDT model；
    \item 调用
      \begin{itemize}
        \item \texttt{solver.boundstatesinrange([...])} 找该 \((L,F)\) manifold 里的束缚态；
        \item \texttt{solver.channelcontributions(nub)} 计算 MQDT channel mixing；
        \item 对允许的 \(m_F=m_i-1,m_i,m_i+1\) 构造 \texttt{state\_mqdt}。
      \end{itemize}
    \item 因此瓶颈不是最后的求和，而是“先把候选终态全部生成出来”。
  \end{itemize}
\end{frame}

\begin{frame}{一个容易忽略的实现细节}
  \begin{block}{函数签名}
    \texttt{total\_decay(self, st, env, printstates=False, include\_fn=lambda s: True)}
  \end{block}

  \begin{itemize}
    \item 签名里有 \texttt{include\_fn} 参数。
    \item 但当前实现中，传给 \texttt{sb.fill} 的是内部固定 lambda：
      \[
        s.n < st.n + 10.
      \]
    \item 因此外部传入的 \texttt{include\_fn} 在当前代码路径下没有真正生效。
    \item 如果要严格控制终态范围，当前更稳妥的方法是显式构造终态列表，然后逐项调用 \texttt{partial\_decay}。
  \end{itemize}
\end{frame}

\begin{frame}[fragile]{问题一：\texttt{st.n} 的语义变了}
  \small
  输入目标态时，我们写的是：
  \begin{verbatim}
initial_state = yb.get_state((65, 0, 1.5, -1.5), tt="vlfm")
  \end{verbatim}

  \vspace{-0.5em}
  物理上想表达：
  \[
    \ket{r}=\ket{65\,{}^3S_1,F=3/2,m_F=-3/2}.
  \]

  但 \texttt{Ytterbium171.get\_state} 会先求 MQDT 根：
  \begin{verbatim}
nutheor = solver.boundstates(nuexp)
nuapprox = round(nutheor * 100) / 100
st = state_mqdt(..., (nuapprox, (-1)**l, f, m), tt="vpfm")
  \end{verbatim}

  \vspace{-0.5em}
  对 \texttt{tt="vpfm"}，\texttt{state\_mqdt} 设置：
  \[
    \texttt{st.n}=\texttt{qn[0]}=\nu_{\mathrm{approx}},
    \qquad
    \texttt{st.v}=\nu_{\mathrm{approx}}.
  \]
  所以这里的 \texttt{st.n} 已经不是整数主量子数 65。
\end{frame}

\begin{frame}{目标态对象的实际数值}
  本地 notebook 中目标态对象打印为：
  \[
    \texttt{|171Yb:64.56,L=0,F=1.5,-1.5>}.
  \]

  \begin{table}
    \centering
    \small
    \begin{tabular}{ll}
      \toprule
      属性 & 数值 / 含义 \\
      \midrule
      输入标签 & \(65\,{}^3S_1,F=3/2,m_F=-3/2\) \\
      \texttt{st.qn} & \texttt{(64.56, 1, 1.5, -1.5)} \\
      \texttt{st.n} & \(64.56\) \\
      \texttt{st.v} & \(64.56\) \\
      \texttt{st.nub} & \(64.56106363\) \\
      \texttt{st.v\_exact} & \(64.56106363\) \\
      \bottomrule
    \end{tabular}
  \end{table}

  \begin{block}{关键错位}
    通用 \texttt{total\_decay} 把 \texttt{st.n} 当“主量子数 \(n\)”使用；
    但 \({}^{171}\mathrm{Yb}\) MQDT 态里，\texttt{st.n} 是四舍五入后的有效量子数 \(\nu\)。
  \end{block}
\end{frame}

\begin{frame}[fragile]{问题二：\texttt{total\_decay} 把 \texttt{st.n} 当搜索半宽}
  \scriptsize
  \texttt{total\_decay} 中的关键代码是：
  \begin{verbatim}
nmax = st.n + 10
sb.fill(st, {
    'dn': st.n,
    'dl': 1,
    'dm': 1,
    'include_fn': lambda s, s0: True if s.n < nmax else False,
})
  \end{verbatim}

  原始设计意图是：
  \begin{itemize}
    \item 若初态是 \(n=65\)，取 \texttt{dn=65}；
    \item 这样就能“从低 \(n\) 到高 \(n\)”自动扫一遍可能衰变终态；
    \item 再用 \(\texttt{s.n}<\texttt{st.n}+10\) 限制最高终态。
  \end{itemize}

  对普通主量子数模型，这个写法虽然粗，但语义还算一致。
  对 Yb171 MQDT，\texttt{dn=st.n} 实际变成：
  \[
    \Delta\nu = 64.56.
  \]
\end{frame}

\begin{frame}[fragile]{问题三：\texttt{get\_nearby} 直接把 \texttt{dn} 用在 \(\nu\) 上}
  \small
  \({}^{171}\mathrm{Yb}\) 的 \texttt{get\_nearby} 里：
  \begin{verbatim}
nu0 = st.nub
...
boundstatesinrange = solver.boundstatesinrange([
    nu0 - o['dn'],
    nu0 + o['dn'],
])
  \end{verbatim}

  因此对当前态：
  \[
    \nu_0=64.56106363,
    \qquad
    \Delta\nu=\texttt{st.n}=64.56.
  \]

  自动搜索区间变成：
  \[
    [\nu_0-\Delta\nu,\nu_0+\Delta\nu]
    =
    [0.00106363,129.12106363].
  \]

  这就是搜索范围被扩大的直接代码原因：\texttt{total\_decay} 的“从 \(n=0\) 扫到附近态”意图，在 MQDT 表示中变成了“从 \(\nu\simeq0\) 扫到 \(\nu\simeq129\)”。
\end{frame}

\begin{frame}[fragile]{问题四：过滤发生得太晚}
  \small
  \texttt{single\_basis.fill} 的顺序是：
  \begin{verbatim}
st_list = s0.atom.get_nearby(self.s0, include_opts=self.opts)

for newst in st_list:
    dipole_ok = s0.allowed_multipole(newst, k=1)
    fn_ok = self.opts['include_fn'](newst, s0)
    if (fn_ok) and (dipole_ok or (not self.opts['dipole_allowed'])):
        self.add(newst)
  \end{verbatim}

  \begin{itemize}
    \item \texttt{include\_fn} 只在 \texttt{get\_nearby} 返回之后才执行。
    \item 因此即使最终过滤条件是
      \[
        \texttt{s.n < st.n + 10}\approx 74.56,
      \]
      \(\nu>74.56\) 的候选也已经被 MQDT 求根、去重和构造过了。
    \item \texttt{dipole\_allowed} 默认是 \texttt{False}，所以 basis 生成阶段也没有只保留 E1 允许终态。
  \end{itemize}

  结论：过滤条件控制的是“最后加入 basis 的状态”，不是“MQDT 求根的搜索范围”。
\end{frame}

\begin{frame}[fragile]{问题五：MQDT 求根本身按细网格扫区间}
  \scriptsize
  \texttt{boundstatesinrange} 的核心循环是：
  \begin{verbatim}
for nubguess in np.arange(range[0], range[1], 0.01):
    nutheorb.append(np.round(self.boundstates(nubguess), decimals=accuracy))
...
nutheorb = [
    item for item in nutheorb
    if channelcontributions(item) passes the l <= nu_i filter
]
  \end{verbatim}

  \begin{itemize}
    \item 区间长度越大，\texttt{boundstates} 调用次数线性增加。
    \item 原生区间长度约为：
      \[
        129.12106363-0.00106363 \simeq 129.12,
      \]
      步长 \(0.01\)，也就是每个 MQDT model 约 \(1.3\times10^4\) 个初始点。
    \item 后续还要对候选根反复调用 \texttt{channelcontributions}。
    \item 低 \(\nu\) 区域会触发 \texttt{sqrt} / fractional power 的非有限数值问题。
  \end{itemize}
\end{frame}

\begin{frame}{本地重跑的诊断结果}
  只测试候选生成，不完整求和，原生 \texttt{total\_decay} 搜索区间为：
  \[
    [0.00106363,129.12106363].
  \]

  \begin{table}
    \centering
    \scriptsize
    \begin{tabular}{lrrl}
      \toprule
      manifold & 候选根数 & 用时 & 备注 \\
      \midrule
      \(L=0,F=1/2\) & -- & \(>25\) s & 单个 solver 超时 \\
      \(L=0,F=3/2\) & 1041 & 0.09 s & 出现极端非物理 \(\nu\) \\
      \(L=1,F=1/2\) & 395 & 10.99 s & P 态通道 \\
      \(L=1,F=3/2\) & 563 & 18.26 s & P 态通道 \\
      \(L=1,F=5/2\) & 516 & 19.96 s & P 态通道 \\
      \bottomrule
    \end{tabular}
  \end{table}

  \begin{itemize}
    \item 这还没有进入所有终态的 \texttt{partial\_decay} 求和。
    \item notebook 中直接调用 \texttt{yb.total\_decay(initial\_state, env)} 在 0 K 和 300 K 都 60 s 超时。
  \end{itemize}
\end{frame}

\begin{frame}{问题本质}
  \begin{block}{不是 \texttt{partial\_decay} 公式失败}
    给定一个初态和一个终态，\texttt{partial\_decay} 仍然可以正常使用电偶极速率公式计算 partial rate。
  \end{block}

  \begin{block}{真正的问题是自动终态集合}
    \texttt{total\_decay} 的自动 basis 生成逻辑来自更通用的氢样 / 碱金属模型；
    它把 \texttt{st.n} 当作主量子数和搜索半宽。
    在 \({}^{171}\mathrm{Yb}\) MQDT 态中，\texttt{st.n} 实际是有效量子数近似值。
  \end{block}

  \begin{itemize}
    \item 语义错位导致搜索区间从 \(\nu\simeq0\) 到 \(\nu\simeq129\)。
    \item 过滤条件在 MQDT 求根之后才生效，不能减少前端成本。
    \item 低 \(\nu\) MQDT 区域数值病态，容易出现非有限系数或非物理解。
  \end{itemize}
\end{frame}

\begin{frame}{第一版处理方式：显式 P 态 subset}
  初始 notebook 没有继续依赖原生 \texttt{total\_decay}，而是显式构造候选终态：
  \begin{itemize}
    \item 第一组诊断：只扫 \(P\) 态 MQDT subset，
      \[
        \nu_{\min}=3.0,\qquad \nu_{\max}=\nu_i+10.
      \]
    \item 第二组诊断：把低 \(n\,{}^3P_J\) 的 NIST 态和高 \(\nu\) MQDT \(P\) 态合并；
    \item 对每个候选终态单独调用：
      \[
        \Gamma_{i\rightarrow f}
        =
        \texttt{partial\_decay}(i,f,\mathrm{env});
      \]
    \item 最后按 partial rate 排序，检查主要贡献通道和截断误差。
  \end{itemize}

  这种方法换来了可控的终态范围和可诊断的主要 decay channels；
  但用户指出寿命明显偏长，说明第一版仍漏掉了重要 channel。
\end{frame}

\begin{frame}{在 rydcalc 基础上做的改进}
  后续 notebook 没有修改 \texttt{rydcalc} 子模块源码，而是在调用层补了三层逻辑：
  \begin{enumerate}
    \item \textbf{分窗 \texttt{get\_nearby}}：
      保留 \texttt{total\_decay} 的“动态生成终态”思想，但把 MQDT \(\nu\) 区间切成小窗口。
    \item \textbf{NIST 低态补充}：
      把 \texttt{Yb171\_NIST.txt} / \texttt{Yb174\_NIST.txt} 中的低 lying E1 终态加入候选池。
    \item \textbf{residual sink}：
      对文献有总寿命的情形，把有限 channel 仍未解释的 rate 作为有效损失通道单独报告。
  \end{enumerate}

  \begin{block}{原则}
    \[
      \Gamma_{\mathrm{model}}
      =
      \Gamma_{\mathrm{NIST}}
      +
      \Gamma_{\mathrm{MQDT}}
      +
      \Gamma_{\mathrm{res}}.
    \]
    其中 \(\Gamma_{\mathrm{res}}\) 只用于总损失校准，不当作真实 resolved final state。
  \end{block}
\end{frame}

\begin{frame}[fragile]{改进一：分窗 \texttt{get\_nearby}}
  \small
  原生 \texttt{total\_decay} 一次请求：
  \[
    \nu\in[0.00106363,129.12106363].
  \]

  notebook-local 修复是构造一个浅拷贝 \texttt{proxy\_state}：
  \begin{verbatim}
proxy.nub = window_center
atom.get_nearby(proxy, {
    "dn": window_half_width,
    "dl": 1,
    "df": 1,
    "dm": 1,
})
  \end{verbatim}

  \begin{itemize}
    \item 仍然由 \texttt{get\_nearby} 动态找 MQDT 终态；
    \item 不再一次扫到极低 \(\nu\) 病态区域；
    \item 相邻窗口有 overlap，再用 \((E,F,m_F,\mathrm{parity})\) 去重；
    \item 对 65S 目标态，窗口法可稳定生成 \(1680\) 个 MQDT P 态候选。
  \end{itemize}
\end{frame}

\begin{frame}{改进二：NIST 低态补充与去重}
  第一版只靠 MQDT \(P\) 态时，低 lying 态贡献不足或能级不够可靠。

  \begin{itemize}
    \item 对 \({}^{171}\mathrm{Yb}\) 65S：
      \begin{itemize}
        \item 显式加入 \texttt{Yb171\_NIST.txt} 中低 \(nP\) 态；
        \item 对每个低 \(P\) 态枚举允许的 \(F,m_F\)；
        \item 使用 \texttt{allowed\_multipole(..., k=1)} 保留 E1 通道；
        \item 高 Rydberg \(P\) 态使用 MQDT，实际取 \(\nu\ge 12\)。
      \end{itemize}
    \item 对 \({}^{174}\mathrm{Yb}\) Fang benchmark：
      \begin{itemize}
        \item 使用 \texttt{Yb174\_NIST.txt} 中的 singlet 低态；
        \item 与 windowed MQDT 高态合并。
      \end{itemize}
  \end{itemize}

  \begin{block}{去重策略}
    相同或近似相同终态按 \((E,F,m,\mathrm{parity})\) 合并；
    低能区优先使用 NIST 能级，高能 Rydberg 区使用 MQDT 态。
  \end{block}
\end{frame}

\begin{frame}{改进三：residual sink}
  有限候选池即使加入 NIST 低态，仍可能少算：
  \begin{itemize}
    \item 其它未解析 manifold；
    \item perturber / configuration-interaction 相关通道；
    \item 实验条件下的 trap-induced loss 或 ionization；
    \item 当前模型没有显式表示的黑体诱导转移。
  \end{itemize}

  因此如果有文献总寿命 \(\tau_{\mathrm{ref}}\)，定义：
  \[
    \Gamma_{\mathrm{ref}}=\frac{1}{\tau_{\mathrm{ref}}},
    \qquad
    \Gamma_{\mathrm{res}}
    =
    \max\left(0,\Gamma_{\mathrm{ref}}-\Gamma_{\mathrm{explicit}}\right).
  \]

  \begin{block}{在动力学模型中的含义}
    \[
      L_{\mathrm{res}}=\sqrt{\Gamma_{\mathrm{res}}}\ket{\mathrm{sink}}\bra{r}.
    \]
    它只保证总人口损失率正确，不提供真实 branching ratio。
  \end{block}
\end{frame}

\begin{frame}{partial\_decay 计算的物理量}
  对一个初态 \(i\) 和终态 \(f\)，先计算角频率：
  \[
    \omega = 2\pi(E_i-E_f).
  \]

  \begin{itemize}
    \item \(\omega>0\)：向低能态辐射，包含自发辐射与受激辐射因子 \(1+n_{\mathrm{th}}\)。
    \item \(\omega<0\)：向高能态吸收黑体光子，只有热光子因子 \(n_{\mathrm{th}}\)。
    \item \(\omega=0\)：代码用一个很大的占据数兜底，正常物理通道中应避免这种情况。
  \end{itemize}

  黑体占据数：
  \[
    n_{\mathrm{th}}(\omega,T)
    =
    \frac{1}{\exp\left(\frac{\hbar|\omega|}{k_B T}\right)-1}.
  \]
\end{frame}

\begin{frame}{partial\_decay 的速率公式}
  电偶极矩阵元由
  \[
    d_{if}=\texttt{st.get\_multipole\_me(other,k=1)}
  \]
  给出，单位是 \(ea_0\)。代码中的速率为：

  \[
    \Gamma_{i\rightarrow f}
    =
    \rho_{\mathrm{LDOS}}\left(\frac{|\omega|}{2\pi}, q\right)
    \cdot
    \eta(\omega,T)
    \cdot
    |\omega|^3
    |d_{if}|^2
    \frac{(e a_0)^2}{3\pi\epsilon_0\hbar c^3}.
  \]

  其中
  \[
    q=m_f-m_i,\qquad
    \eta(\omega,T)=
    \begin{cases}
      1+n_{\mathrm{th}}, & \omega>0,\\
      n_{\mathrm{th}}, & \omega<0.
    \end{cases}
  \]
\end{frame}

\begin{frame}{矩阵元内部如何处理 MQDT 态}
  \texttt{get\_multipole\_me} 对 MQDT 态不是简单查表，而是对 channel 逐项求和：
  \[
    \ket{i}=\sum_\alpha A_\alpha^{(i)}\ket{\alpha},
    \qquad
    \ket{f}=\sum_\beta A_\beta^{(f)}\ket{\beta}.
  \]

  对每对 channel：
  \begin{itemize}
    \item 检查 core state 是否相容；
    \item 用 Wigner \(3j/6j\) 系数处理角动量约化矩阵元；
    \item 用径向积分 \(\langle R_f|r^k|R_i\rangle\) 给出 radial part；
    \item 对所有 channel pair 加权求和。
  \end{itemize}

  因此寿命结果同时依赖能级、角动量代数、径向波函数和 MQDT channel mixing。
\end{frame}

\begin{frame}{total\_decay 与显式终态求和的关系}
  两者使用相同的 partial rate 公式：
  \[
    \Gamma_{\mathrm{tot}}=\sum_f\Gamma_{i\rightarrow f}.
  \]

  差别只在“终态集合”：
  \begin{itemize}
    \item \texttt{total\_decay}：内部自动调用 \texttt{single\_basis.fill}，终态范围由当前实现固定。
    \item 显式求和：用户自己构造 \(\{f\}\)，再逐项调用 \texttt{partial\_decay}。
  \end{itemize}

  \begin{block}{什么时候应该显式构造终态}
    当自动终态搜索过大、过慢、扫到数值病态区，或者需要明确报告截断条件时，显式构造终态更可控。
    但这只解决“可控性”，不自动保证“完备性”；漏掉低能或其它 manifold 会直接导致寿命偏长。
  \end{block}
\end{frame}

\begin{frame}{对 \({}^{171}\mathrm{Yb}\) 寿命计算的经验}
  以
  \[
    \ket{r}=\ket{65\,{}^3S_1,F=3/2,m_F=-3/2}
  \]
  为例：

  \begin{itemize}
    \item 态构造应使用：
      \[
        \texttt{get\_state((65, 0, 1.5, -1.5), tt="vlfm")}.
      \]
    \item 原生 \texttt{total\_decay} 会从很低的有效量子数区域开始自动生成终态。
    \item 在当前本地环境中，直接调用曾出现长时间不返回和低 \(\nu\) 插值 / 非有限系数问题。
    \item notebook 中采用的显式 P 态终态范围：
      \[
        \nu_{\min}=3.0,\qquad
        \nu_{\max}=\nu_i+10.
      \]
    \item 这个范围会漏掉低 \(\nu\) P 态，因此只能作为诊断性 subset，不能作为完整寿命。
  \end{itemize}
\end{frame}

\begin{frame}{早期结果：65 \({}^3S_1\) 态的诊断性 subset}
  第一版 \(\nu_{\min}=3.0\) MQDT-only 终态求和结果：
  \begin{table}
    \centering
    \begin{tabular}{lrrr}
      \toprule
      温度 & 候选终态数 & \(\Gamma_{\mathrm{tot}}\) / s\(^{-1}\) & \(\tau\) / \(\mu\)s \\
      \midrule
      0 K & 1680 & \(1.554823\times 10^3\) & 643.160 \\
      300 K & 1680 & \(5.596472\times 10^3\) & 178.684 \\
      \bottomrule
    \end{tabular}
  \end{table}

  \begin{itemize}
    \item 这些寿命明显偏长，说明这个终态集合少算了 channel。
    \item 后续改进的目标不是改变 \texttt{partial\_decay} 公式，而是补全终态集合。
    \item 这个表应理解为“缺 channel 前”的 baseline。
  \end{itemize}
\end{frame}

\begin{frame}{Fang 2001 benchmark：低 \(n\) Yb 寿命}
  用 \({}^{174}\mathrm{Yb}\) 的 \(6sns\,{}^1S_0\) 和 \(6snd\,{}^1D_2\) 序列做验证。
  Fang et al. 给出了室温 \(\tau=1/(A+B\rho)\) 计算值。

  \begin{table}
    \centering
    \scriptsize
    \begin{tabular}{lrrrr}
      \toprule
      态 & hybrid 显式 \(\tau_{300K}\) & Fang \(\tau_{300K}\) & \(\Gamma_{\mathrm{res}}\) / s\(^{-1}\) & 校准后 \(\tau\) \\
      \midrule
      \(6s21s\) & 7.424 & 6.8 & \(1.235\times10^4\) & 6.800 \\
      \(6s22s\) & 8.417 & 8.3 & \(1.671\times10^3\) & 8.300 \\
      \(6s27s\) & 14.508 & 16.9 & 0 & 14.508 \\
      \(6s21d\) & 15.218 & 7.8 & \(6.249\times10^4\) & 7.800 \\
      \(6s27d\) & 28.090 & 18.2 & \(1.934\times10^4\) & 18.200 \\
      \bottomrule
    \end{tabular}
  \end{table}

  \begin{itemize}
    \item \(s\) 系列在加入 NIST 低态后已接近 Fang 的量级和趋势；
    \item \(d\) 系列仍需要较大的 residual，说明低 \(n\) perturbation / configuration mixing 仍未被显式模型完整覆盖。
  \end{itemize}
\end{frame}

\begin{frame}{最终目标态：explicit hybrid 结果}
  对
  \[
    \ket{r}=\ket{65\,{}^3S_1,F=3/2,m_F=-3/2}
  \]
  使用 NIST low \(P\) states + MQDT \(P\) states with \(\nu\ge12\)：

  \begin{table}
    \centering
    \begin{tabular}{lrr}
      \toprule
      项目 & 0 K & 300 K \\
      \midrule
      候选终态数 & \multicolumn{2}{c}{1542 = 84 NIST + 1458 MQDT} \\
      \(\Gamma_{\mathrm{NIST}}\) / s\(^{-1}\) & -- & \(3.041537\times10^3\) \\
      \(\Gamma_{\mathrm{MQDT}}\) / s\(^{-1}\) & -- & \(4.261963\times10^3\) \\
      \(\Gamma_{\mathrm{explicit}}\) / s\(^{-1}\) & \(3.262577\times10^3\) & \(7.303500\times10^3\) \\
      \(\tau_{\mathrm{explicit}}\) / \(\mu\)s & 306.506 & 136.921 \\
      \bottomrule
    \end{tabular}
  \end{table}

  \begin{itemize}
    \item 相比 MQDT-only subset，300 K 显式速率从 \(5.60\times10^3\) 升到 \(7.30\times10^3\,\mathrm{s}^{-1}\)。
    \item 寿命仍长于实验值，说明还存在未解析损失。
  \end{itemize}
\end{frame}

\begin{frame}{目标态：与实验寿命的 residual 校准}
  Muniz et al. 对相同 Rydberg 态在 typical tweezer depth 下给出：
  \[
    \tau_{\mathrm{exp}}\simeq65(3)\,\mu\mathrm{s}.
  \]

  因此：
  \[
    \Gamma_{\mathrm{ref}}
    =
    \frac{1}{65\,\mu\mathrm{s}}
    =
    1.538462\times10^4\,\mathrm{s}^{-1}.
  \]

  \begin{table}
    \centering
    \begin{tabular}{lr}
      \toprule
      贡献 & rate / s\(^{-1}\) \\
      \midrule
      explicit hybrid & \(7.303500\times10^3\) \\
      residual sink & \(8.081115\times10^3\) \\
      calibrated total & \(1.538462\times10^4\) \\
      \bottomrule
    \end{tabular}
  \end{table}

  \begin{block}{使用建议}
    做 branching 分析时看 explicit NIST/MQDT channels；
    做 gate loss / erasure 模型时再加入 residual sink，使总寿命匹配实验。
  \end{block}
\end{frame}

\begin{frame}{主要 300 K explicit 通道}
  \begin{table}
    \centering
    \scriptsize
    \begin{tabular}{rllr}
      \toprule
      排名 & 来源 & 终态 & \(\Gamma\) / s\(^{-1}\) \\
      \midrule
      1 & NIST & \(6\,{}^3P_2,F=5/2,m_F=-5/2\) & 544.428 \\
      2 & MQDT & \(\nu=65.08,L=1,F=5/2,m_F=-5/2\) & 510.435 \\
      3 & MQDT & \(\nu=64.08,L=1,F=5/2,m_F=-5/2\) & 483.589 \\
      4 & NIST & \(6\,{}^3P_1,F=3/2,m_F=-3/2\) & 277.815 \\
      5 & MQDT & \(\nu=65.27,L=1,F=3/2,m_F=-3/2\) & 256.884 \\
      6 & MQDT & \(\nu=65.34,L=1,F=1/2,m_F=-1/2\) & 218.292 \\
      \bottomrule
    \end{tabular}
  \end{table}

  \begin{itemize}
    \item NIST low \(P\) 态负责大能量差的 spontaneous decay；
    \item 相邻 Rydberg \(P\) 态负责 300 K 下显著的 BBR-induced redistribution；
    \item 这两类通道必须同时保留，才能接近真实 lifetime scale。
  \end{itemize}
\end{frame}

\begin{frame}{实用建议}
  \begin{enumerate}
    \item 先确认态标记格式：Yb171 Rydberg MQDT 态优先用 \texttt{vlfm}。
    \item 记录环境参数：至少写明 \(T\)、LDOS 是否为默认值。
    \item 对 Yb171 MQDT 态，不直接依赖原生 \texttt{total\_decay} 的一次性大范围搜索。
    \item 优先使用 windowed \texttt{get\_nearby} 构造高 Rydberg 终态，并显式报告 \(\nu\) 范围。
    \item 对低 lying final states，优先使用 NIST 能级并与 MQDT 候选去重。
    \item 如果要匹配实验总寿命，把未解析部分作为 residual sink 单独列出。
    \item 发表或 PR 中报告寿命时，同时报告：
      \begin{itemize}
        \item 初态；
        \item 终态截断；
        \item 是否包含 BBR；
        \item \(\Gamma\) 与 \(\tau\)；
        \item 主要 partial decay channels。
      \end{itemize}
  \end{enumerate}
\end{frame}

\begin{frame}{总结}
  \begin{itemize}
    \item \texttt{rydcalc} 提供从原子结构到相互作用、矩阵元和寿命的统一计算接口。
    \item \texttt{total\_decay} 的本质是：
      \[
        \text{自动生成终态}\quad+\quad
        \sum_f \texttt{partial\_decay}(i,f).
      \]
    \item \texttt{partial\_decay} 使用标准电偶极辐射速率公式，并加入温度相关的黑体光子占据数。
    \item 对复杂 MQDT 体系，寿命计算的可靠性不只取决于公式，还取决于终态集合是否数值稳定且物理完备。
    \item 我们在不改 \texttt{rydcalc} 源码的前提下，引入了分窗 \texttt{get\_nearby}、NIST 低态补充和 residual sink 校准。
    \item 对当前 \({}^{171}\mathrm{Yb}\) 65S 目标态，explicit hybrid 给出
      \(\tau_{300K}=136.921\,\mu\mathrm{s}\)；
      加 residual sink 后可匹配 \(65\,\mu\mathrm{s}\) 实验总寿命。
  \end{itemize}
\end{frame}

\end{document}
